% Working evaluation-section fragment, 2026-10-31.
% Protocol intent, not evidence that all Core100 gates or experiments passed.
% Table counts must be filled from versioned generation/admission/label receipts.
% Current evidence: review/cvpr_experiments_20261031/CORE100_EVALUATION_CHECKLIST.json
% Input comparison: review/cvpr_experiments_20261031/experiment_comparability/REPORT.json

\subsection{Evaluation Protocol}
\label{sec:evaluation_protocol}

\paragraph{Physical response at three levels.}
The evaluation asks whether a generated continuation responds to a declared
physical intervention. We examine the resulting event branch, the downstream
event sequence, and quantitative dynamics where measurement is identifiable.
These are complementary observations of the same intervention-response
problem, with separate applicability conditions. They are reported
separately rather than combined through arbitrary weights into one score.

\paragraph{Inputs and evaluation units.}
The test membership consists of 100 physical cases, split into ID80,
OOD-1 (T03/T20, ten cases), and OOD-2 (T11/T19, ten cases). Each case has three
English prompt variants. The primary generation setting uses seed zero for
each variant. A generation slot is identified by model revision, input
release, case, variant, seed, and sampling configuration; the actual prompt,
condition image, and output video are bound by hashes. Three prompt variants
are repeated conditions on one physical case, not three random seeds or
independent trajectories.

\paragraph{Reference and window admission.}
Each target has a versioned event definition, a registered time origin and
observation horizon, and a corresponding source-state record. Boundary
evidence must use the same event, fixed physical context, intervention axis,
input version, and observation window. A parameter-scan transition is not
automatically a unique critical point: multiple observed transitions and
unresolved interval distances are retained explicitly. Published scan points
with different simulation seeds are distinguished from fixed-source-seed
interventions. Source-event reproduction, RGB observability, and the
numerical measurements required by a predicate are checked separately.

For each generated output, native frame count, frame rate, decoding, and
window coverage are verified. A shorter video cannot establish that an event
does not occur later in the registered window. Playback is not slowed, frames
are not repeated, and a static lead-in is not silently removed to make a
video satisfy the protocol. Short-window observations require their own
matching reference and boundary evidence before they can support a
short-window physical score.

\paragraph{VLM-assisted human assessment.}
A frozen visual-language-model pipeline supplies preannotations for scene
integrity and supported visible events. Scene integrity covers newly
appearing foreign objects, scene replacement, loss of original-object
identity, and severe geometry corruption. Such machine flags are review
candidates; all registered videos remain in the inventory. Human reviewers
can assess a video without waiting for a machine prediction. Disagreements
and unobservable events remain explicit until resolved.

Human event observations do not require invented timestamps when the
registered form asks only for event labels. Where a sequence requires an
order judgement, separate confirmation of the event order is required;
two positive event labels do not establish their order. Likewise, an RGB
judgement is not substituted for a calibrated displacement, angle, velocity,
or other numerical predicate. An independent evaluation of the automatic
assessor requires fresh human observations with machine predictions hidden;
the production VLM-assisted labels are not described as independent gold.

\paragraph{Physical metrics and coverage.}
Branch accuracy compares admitted target-event observations with the
corresponding reference branch. NearAcc evaluates the registered near-boundary
subset using the frozen distance protocol. PairAcc requires both admitted
opposite-branch endpoints of a registered pair to be correct under the same
prompt variant and coupled seed; it is not the product of two marginal
accuracies. EventExact requires the registered event bits and required order,
whereas event-bit agreement alone is reported with a distinct label.
GravTrace provides a separate quantitative gravity diagnostic on videos with
the required trajectory, projection, clock, and identifiability evidence.

Each table must identify the eligible reference set, generated-video
coverage, unresolved annotations, measurement coverage, and confirmed
generated failures. Missing reference evidence or an unreviewed label is not
assigned an observed zero. A confirmed generated failure on an admitted task
remains in that task's denominator under the frozen failure policy; exclusion
by a machine quality flag alone is not permitted.

\paragraph{Training-utility comparisons.}
Base and adapted models are compared within a matched experiment using the
same base revision, input manifest, native frame grid, resolution, sampler,
seed, and prompt variant. Fixed training checkpoints are selected by the
declared schedule rather than test performance. Actual optimizer steps,
global batch size, and video presentations must be reported from runtime
receipts; equal step numbers alone do not prove equal training budgets.
Outputs from a different prompt release or native generation regime are not
substituted into the original multi-model table.

\paragraph{Statistical units.}
Results are separated by prompt variant and ID/OOD split. Any aggregate
retains the registered physical source groups as the uncertainty unit,
keeping synchronized views, verified relatives, and repeated prompt
conditions together. Small OOD sets are accompanied by counts and coverage;
their prompt repetitions are not treated as independent physical trials.

% Pending before final submission:
% - independent human assessor validation and adjudication;
% - current-reference/full-window admission and generated event labels;
% - calibrated quantitative measurements and source-group uncertainty results;
% - completed matched training-utility outputs and fixed-budget receipts.
